% Canonical continuation of cusp_spectral_geometry.tex.  It uses the notation
% and labels of cusp_globalization_replacement.tex and
% cusp_cone_sphere_routes.tex, which are included immediately before it.

\subsubsection{Optimal sphere exponents and the exact Euclidean boundary}
\label{cusp-cone-optimal-strengthening}

Within the class of native connected fibrewise torus quotients equipped with
the induced, unscaled quotient of the marked Euclidean metric, the two sphere
dimensions in the preceding constructions are forced by two different
integral quotient problems.  Cusp-locally, each copy has an expanding
quotient of rank two, so two copies are necessary and sufficient within that
class.  Over the full regular base, monodromy leaves at most one expanding
quotient direction per copy, so four copies are necessary and sufficient
within the same class.  We now prove this scoped optimality, identify the
compact conical carriers, and construct a noncompact comparison manifold in
which the limiting \(\mathbb R^4\) is an actual boundary fibre.  There is only
a continuous comparison map from this manifold to the compact conical
carrier: it set-theoretically collapses the new boundary, is a homeomorphism
away from that boundary, and is nonproper and nonquotient.

Throughout this subsubsection,
\[
 B_R^Z(z):=\{w\in Z:d_Z(z,w)\leq R\}
\]
denotes the closed radius-\(R\) ball; when the ambient space and center are
clear, superscripts and centers are suppressed.

\begin{theorem}[Optimality within the native connected torus-quotient class]
\label{cusp-cone-optimal-exponents}
Consider only native connected fibrewise torus quotients of copies of the
original marked fibres, and equip every such quotient with the induced,
unscaled Euclidean quotient metric.

\begin{enumerate}
\item For a native connected torus quotient of \(r\) copies over the
punctured cusp, if the quotient systole tends to infinity on the positive
radial ray, then its lattice rank is at most \(2r\).  A rank-four quotient
therefore requires \(r\geq2\).  Equality is attained for \(r=2\) by
\[
 (\Lambda/K_\Lambda)^{\oplus2},
 \qquad K_\Lambda=
 \mathbb Z\langle\widehat w,\widehat\delta\rangle,
\]
whose physical period lattice is
\[
 L_T^{(12)}=\operatorname{diag}(Y_T,Y_T)\mathbb Z^4.
\]
The corresponding smash source is
\(X^{\wedge2}\cong S^{12}\).

\item For a native connected torus quotient of \(r\) copies over the full
regular base, or over a finite cover of that base, radial divergence of the
quotient systole forces its lattice rank to be at most \(r\).  A rank-four
quotient therefore requires \(r\geq4\).  Equality is attained by the four
copies of the full-monodromy-invariant character \(\chi_\gamma\); on the
radial fibre its period lattice is
\[
 L_T^{(24)}=\ell_T\mathbb Z^4,
 \qquad
 \ell_T=T+H-B+O(T^{-1}),
\]
and the corresponding smash source is
\(X^{\wedge4}\cong S^{24}\).
\end{enumerate}

Thus, among precisely these native connected torus quotients with their
induced unscaled marked Euclidean metrics, \(S^{12}\) is the smallest smash
sphere produced by the cusp-local rank-two quotient, and \(S^{24}\) is the
smallest smash sphere produced by a quotient which survives the full
regular-base monodromy.  No optimality is asserted here among arbitrary
continuous quotient maps or after a \(T\)-dependent rescaling of the metric.
\end{theorem}

\begin{proof}
For \(r\) copies, every element of
\[
 K_\Lambda^{\oplus r}
 =\mathbb Z\langle\widehat w,\widehat\delta\rangle^{\oplus r}
\]
has uniformly bounded physical length on the radial cusp.  Let
\(N\subset\Lambda^{\oplus r}\) be a primitive sublattice defining a
connected torus quotient.  If some \(v\in K_\Lambda^{\oplus r}\) had
nonzero class in \(\Lambda^{\oplus r}/N\), that nonzero quotient period
would have quotient norm at most the original norm of \(v\), uniformly in
\(T\).  The quotient systole could not diverge.  Hence
\[
 K_\Lambda^{\oplus r}\subset N,
\]
and
\[
 \operatorname{rank}(\Lambda^{\oplus r}/N)
 \leq 4r-2r=2r.
\]
For \(r=2\), the primitive kernel \(K_\Lambda^{\oplus2}\) attains this
bound.  The cusp monodromy is the identity modulo \(K_\Lambda\), so the
quotient descends over the punctured cusp.  Its radial lattice is exactly
\(\operatorname{diag}(Y_T,Y_T)\mathbb Z^4\), and
\[
 \operatorname{sys}(L_T^{(12)})
 =\operatorname{sys}(Y_T\mathbb Z^2)\geq T-C\longrightarrow\infty.
\]

For a quotient over the full regular base, monodromy invariance gives a
stronger bound.  Since \(\widehat w\in N\) and
\[
 A_1\widehat w=\widehat u-\widehat w,
\]
one has \(\widehat u\in N\).  Together with
\(\widehat w,\widehat\delta\in N\), this places the rank-three lattice
\[
 \mathbb Z\langle\widehat u,\widehat w,
 \widehat\delta\rangle^{\oplus r}
\]
inside \(N\), and hence
\[
 \operatorname{rank}(\Lambda^{\oplus r}/N)\leq r.
\]
The same conclusion holds after a finite base cover.  Indeed, if
\(\Gamma_{\mathrm{mon}}\) is the corresponding finite-index monodromy
subgroup and \(G_*=A_2A_1=I+N_*\), then \(N_*^2=0\), and
\(\Gamma_{\mathrm{mon}}\) contains \(G_*^m\) for
some \(m>0\).  The exact orbit formula
\[
 G_*^m\widehat w
 =\widehat w+m(\widehat u-\widehat\delta)
\]
implies \(m\widehat u\in N\); primitivity of \(N\) implies
\(\widehat u\in N\).

The invariant character \(\chi_\gamma\) has precisely the complementary
rank-one quotient in each copy.  The full-base construction, including the
two finite Seifert fillings and the cone at the unipotent cusp, proves that
four copies attain rank four.  Finally, after choosing the based
diffeomorphism \(X\cong_{\mathrm{diff}}S^6\) already reconstructed in this
workbench,
\(X^{\wedge r}\cong(S^6)^{\wedge r}\cong S^{6r}\), which gives the two
claimed sphere dimensions.
\end{proof}

\begin{theorem}[Exact sphere-to-sphere carriers with moving four-spheres]
\label{cusp-cone-angle-forgotten-sphere-carriers}
There are continuous surjective based quotient maps
\[
 \widehat q_3:(X,x_0)\longrightarrow(S^3,*),
 \qquad
 \widehat q_\gamma:(X,x_0)\longrightarrow(S^2,*),
\]
whose smash powers are continuous surjective quotient maps
\begin{equation}
 \widehat q_3^{\wedge2}:S^{12}\twoheadrightarrow S^6,
 \qquad
 \widehat q_\gamma^{\wedge4}:S^{24}\twoheadrightarrow S^8.
 \label{cusp-cone-sphere-to-sphere-maps}
\end{equation}
Use the marked presentation \(\overline T^2=\mathbb R^2/\mathbb Z^2\) and
put
\[
 K^{(1)}
 =((\mathbb R/\mathbb Z)\times\{0\})
  \cup(\{0\}\times(\mathbb R/\mathbb Z)),
 \qquad
 c:\overline T^2\longrightarrow\overline T^2/K^{(1)}\cong S^2.
\]
Here \(c\) is the quotient collapsing the coordinate one-skeleton.  For
every sufficiently large \(T\), put
\begin{align*}
 r_T&=e^{-2\pi T},\\
 K_T^{(1)}&:=\iota_T(K^{(1)})\subset Q_T,\\
 S_T^2&:=\overline T^2/K^{(1)}\cong Q_T/K_T^{(1)},\\
 c_T&:=c\circ\iota_T^{-1}:Q_T\longrightarrow S_T^2,
\end{align*}
where the identification of the two displayed quotients is induced by
\(\iota_T\).  Within this theorem and its proof, \(c_T\) denotes the
displayed collapse map; the earlier scalar \(c_T=\xi_T-\eta_T\) does not
enter.
At reduced-suspension cross-section height \(r_T\), the first target contains
a based topologically embedded four-sphere, equipped with its intrinsic
quotient length metric,
\[
 \mathbb S_{12,T}^4=S_T^2\wedge S_T^2\hookrightarrow S^6,
\]
and the second target contains, at the same cross-section height, a based
topologically embedded four-sphere, equipped with its intrinsic quotient
length metric,
\[
 \mathbb S_{24,T}^4=C_T^{\wedge4}\hookrightarrow S^8.
\]
On synchronized radial fibres, the exact restrictions factor as
\begin{align}
 F_T^2&\xrightarrow{\ \pi_T\times\pi_T\ }
 Q_T^2\xrightarrow{\ c_T\times c_T\ }(S_T^2)^2
 \xrightarrow{\ \rho_{12,T}\ }S_T^2\wedge S_T^2
 =\mathbb S_{12,T}^4,
 \label{cusp-cone-S12-S4-radial-map}\\
 F_T^4&\xrightarrow{\ \chi_{\gamma,T}^{\times4}\ }
 C_T^4\xrightarrow{\ \rho_{24,T}\ }C_T^{\wedge4}
 =\mathbb S_{24,T}^4,
 \label{cusp-cone-S24-S4-radial-map}
\end{align}
Here \(c_T\) collapses \(K_T^{(1)}\), and \(\rho_{12,T}\) and
\(\rho_{24,T}\) are the smash quotient maps.  Thus the
full collapsed locus in \(Q_T^2\) is exactly
\begin{equation}
 W_{12,T}
 =\bigl(K_T^{(1)}\times Q_T\bigr)
  \cup\bigl(Q_T\times K_T^{(1)}\bigr),
 \label{cusp-cone-S12-exact-collapsed-locus}
\end{equation}
and the full collapsed locus in \(C_T^4\) is exactly
\begin{equation}
 W_{24,T}
 =\{(x_1,x_2,x_3,x_4)\in C_T^4:
       x_j=0\pmod{\ell_T}\ \text{for some }j\}.
 \label{cusp-cone-S24-exact-collapsed-locus}
\end{equation}
The maps
\(\rho_{12,T}\circ(c_T\times c_T):Q_T^2\to
S_T^2\wedge S_T^2\) and
\(\rho_{24,T}:C_T^4\to C_T^{\wedge4}\) send exactly \(W_{12,T}\) and
\(W_{24,T}\), respectively, to the smash basepoint and are injective on
their complements.  The full composites from \(F_T^2\) and \(F_T^4\)
still identify every pair of points identified by
\(\pi_T\times\pi_T\) and \(\chi_{\gamma,T}^{\times4}\), respectively; in
particular, their first arrows retain their nontrivial torus kernels.
Their exact smash-basepoint preimages in the full radial sources are
\[
 \bigl(\pi_T^{-1}(K_T^{(1)})\times F_T\bigr)
 \cup\bigl(F_T\times\pi_T^{-1}(K_T^{(1)})\bigr)
\]
and
\[
 \{(x_1,x_2,x_3,x_4)\in F_T^4:
   \chi_{\gamma,T}(x_j)=0\ \text{in }C_T\ \text{for some }j\},
\]
respectively.  Outside these basepoint preimages, two source points have the
same final image if and only if their images under the corresponding first
arrow in \cref{cusp-cone-S12-S4-radial-map,cusp-cone-S24-S4-radial-map}
are equal.

Choose the Euclidean Riemannian \(\ell^2\)-product length metrics on
\(Q_T^2\) and \(C_T^4\) before taking the smash quotients, and use the
successive quotient length metrics specified in the proof.  With the
specified midpoint basepoints,
\begin{equation}
 (\mathbb S_{12,T}^4,\zeta_{12,T})
 \longrightarrow(\mathbb R^4,0),
 \qquad
 (\mathbb S_{24,T}^4,\zeta_{24,T})
 \longrightarrow(\mathbb R^4,0)
 \label{cusp-cone-moving-S4-R4-limits}
\end{equation}
in pointed local-isometric and pointed Gromov--Hausdorff topology.  Moreover,
the convergence is smooth on the eventually Euclidean basepoint
neighborhoods.  For every \(R>0\), once \(T-C>4R\), respectively once
\(\ell_T>4R\), the corresponding quotient map restricts to an isometry from
the radius-\(R\) ball in \(Q_T^2\), respectively \(C_T^4\), \emph{onto} the
entire radius-\(R\) ball in the indicated smash sphere.  The domain ball, and
hence the quotient ball, is exactly a Euclidean four-ball.
\end{theorem}

\begin{proof}
Write \(\Sigma_r(Z,z_0)\) for the reduced suspension
\[
 \frac{[0,\epsilon]\times Z}
 {(\{0,\epsilon\}\times Z)\cup([0,\epsilon]\times\{z_0\})},
\]
where the displayed union is collapsed to one basepoint.

Use the marked presentation
\(\overline T^2=\mathbb R^2/\mathbb Z^2\), let
\[
 K^{(1)}=
 ((\mathbb R/\mathbb Z)\times\{0\})
 \cup(\{0\}\times(\mathbb R/\mathbb Z)),
\]
and let
\[
 c:\overline T^2\longrightarrow
 \overline T^2/K^{(1)}\cong S^2
\]
be the degree-one quotient obtained by collapsing the one-skeleton.  The
formula
\[
 \widehat p_2:\mathcal A\longrightarrow
 \Sigma_r(\overline T^2/K^{(1)},c(K^{(1)}))\cong S^3,
 \qquad
 [r,e^{i\phi},u]\longmapsto[r,c(u)]
\]
is well defined: the two suspension ends of \(\mathcal A\) and the entire
cylinder over \(K^{(1)}\) map to the reduced-suspension basepoint.  It is a
continuous surjection.  Set
\(\widehat q_3=\widehat p_2\circ q_2\).

For every sufficiently large \(T\), define the transported one-skeleton and
its collapse by
\begin{equation}
 \begin{aligned}
 K_T^{(1)}&:=\iota_T(K^{(1)})\subset Q_T,\\
 S_T^2&:=\overline T^2/K^{(1)}\cong Q_T/K_T^{(1)},\\
 c_T&:=c\circ\iota_T^{-1}:Q_T\longrightarrow S_T^2.
 \end{aligned}
 \label{cusp-cone-transported-skeleton-collapse}
\end{equation}
The displayed target identification is the one induced by \(\iota_T\), so
\(c_T^{-1}(*)=K_T^{(1)}\), where
\(*=c(K^{(1)})=c_T(K_T^{(1)})\).

Give \(Q_T\) its Euclidean quotient length metric and \(S_T^2\) the quotient
length metric under \(c_T\).  Put
\[
 y_*=[(1/2,1/2)]\in\overline T^2,
 \qquad y_T:=\iota_T(y_*)\in Q_T,
 \qquad z_T:=c_T(y_T)\in S_T^2.
\]
Equivalently, the metric on \(\overline T^2\) used below is
\(g_T:=\iota_T^*g_{Q_T}\).  For every \(v\in\mathbb R^2\), the exact
decomposition
\[
 Y_T=T\begin{pmatrix}0&1\\-1&0\end{pmatrix}+E_T,
 \qquad\|E_T\|_{\mathrm{op}}\leq C,
\]
gives
\[
 \|Y_Tv\|\geq(T-C)\|v\|.
\]
The standard flat distance from \(y_*\) to \(K^{(1)}\) is \(1/2\);
therefore
\begin{equation}
 d_T:=\operatorname{dist}_{Q_T}(y_T,K_T^{(1)})
 =\operatorname{dist}_{g_T}(y_*,K^{(1)})
 \geq\frac{T-C}{2}.
 \label{cusp-cone-distance-to-skeleton}
\end{equation}
For a length-space quotient which collapses a closed subset \(K\) to one
point, the distance between two noncollapsed classes represented by \(x,y\)
is
\[
 \min\{d(x,y),d(x,K)+d(y,K)\}.
\]
The two displayed terms are upper bounds, obtained respectively from a
direct path and from paths from \(x\) and \(y\) to \(K\).  Conversely, a
rectifiable quotient path either avoids the collapsed point, in which case
its compact image lifts to \(X\setminus K\) and has length at least
\(d(x,y)\), or meets that point, in which case its two pieces have total
length at least \(d(x,K)+d(y,K)\).  Taking infima proves the formula.
If \(T-C>4R\), then \(d_T>2R\).  The displayed formula and the systolic
bound show that
\[
 c_T:B_R^{Q_T}(y_T)\longrightarrow B_R^{S_T^2}(z_T)
\]
is an isometry onto the entire target ball, and both balls are exactly
Euclidean two-balls.  Indeed, the second term in the minimum cannot realize
a distance at most \(2R\), and any point of the target ball has a unique
noncollapsed representative; the same formula puts that representative in
the source ball.

At reduced-suspension coordinate \(r_T=e^{-2\pi T}\), \(S_T^2\) is a based
embedded subspace of \(S^3\): its point \(*\) is the global
reduced-suspension basepoint.  Hence
\[
 \mathbb S_{12,T}^4=S_T^2\wedge S_T^2
\]
embeds in \(S^3\wedge S^3=S^6\).  Equip \(Q_T^2\) first with the Euclidean
Riemannian \(\ell^2\)-product length metric, equip \((S_T^2)^2\) with the
\(\ell^2\)-product of the two quotient length metrics, and let
\[
 \rho_{12,T}:(S_T^2)^2\longrightarrow S_T^2\wedge S_T^2
\]
carry the quotient length metric obtained by collapsing the fat wedge
\[
 \mathcal W_{12,T}
 =(S_T^2\times\{*\})\cup(\{*\}\times S_T^2).
\]
Since \(c_T^{-1}(*)=K_T^{(1)}\), the inverse image of
\(\mathcal W_{12,T}\) under \(c_T\times c_T\) is exactly \(W_{12,T}\) in
\cref{cusp-cone-S12-exact-collapsed-locus}.  On
\(Q_T^2\setminus W_{12,T}\), the map \(c_T\times c_T\) is injective and its
image avoids \(\mathcal W_{12,T}\); the map \(\rho_{12,T}\) is injective on
the complement of \(\mathcal W_{12,T}\).  Hence their composite identifies
all of \(W_{12,T}\) to the smash basepoint and makes no other
identifications.

The point \(\zeta_{12,T}=[z_T,z_T]\) has distance \(d_T\) from this fat
wedge.  The one-factor onto isometry above gives an onto isometry from the
radius-\(R\) ball about \((y_T,y_T)\) in \(Q_T^2\) to the corresponding
ball about \((z_T,z_T)\) in \((S_T^2)^2\).  Applying the same collapsed-set
distance formula to \(\rho_{12,T}\), and using \(d_T>2R\), shows that
\begin{equation}
 \rho_{12,T}\circ(c_T\times c_T):
 B_R^{Q_T^2}((y_T,y_T))
 \xrightarrow{\ \cong_{\mathrm{isom}}\ }
 B_R^{\mathbb S_{12,T}^4}(\zeta_{12,T})
 \label{cusp-cone-S12-onto-exact-ball}
\end{equation}
is onto as well as isometrically embedded.  The systolic bound
\(\operatorname{sys}(Q_T^2)\geq T-C>4R\) makes its domain exactly a
Euclidean four-ball.  On a radial fibre,
\[
 (\widehat p_2\circ q_2)|_{F_T}
 =[r_T,c(\iota_T^{-1}(\pi_T(\,\cdot\,)))]
 =[r_T,c_T(\pi_T(\,\cdot\,))],
\]
and taking two factors followed by \(\rho_{12,T}\) proves the entire
factorization in \cref{cusp-cone-S12-S4-radial-map}.

For the four-copy construction, choose \(1\in S^1\) as character basepoint
and define
\[
 \widehat p_1:\mathcal B\longrightarrow\Sigma_r(S^1,1)\cong S^2,
 \qquad
 [r,e^{i\phi},\theta]\longmapsto[r,\theta].
\]
The formula forgets the base angle and collapses the two suspension ends
and the entire cylinder with \(\theta=1\) to the reduced-suspension
basepoint.  It is a continuous surjection.  Put
\(\widehat q_\gamma=\widehat p_1\circ q_{1,\mathrm{cusp}}\).

At the interior suspension coordinate \(r_T=e^{-2\pi T}\), identify the
character circle \(C_T=\mathbb R/\ell_T\mathbb Z\) with
\(\mathbb R/\mathbb Z\subset S^2\) by \([x]\mapsto[x/\ell_T]\).  It is a
based embedded circle, so its fourth smash is an embedded
\[
 \mathbb S_{24,T}^4=C_T^{\wedge4}\subset(S^2)^{\wedge4}=S^8.
\]
Equip \(C_T^4\) first with the Euclidean Riemannian \(\ell^2\)-product
length metric, and let
\[
 \rho_{24,T}:C_T^4\longrightarrow C_T^{\wedge4}
\]
carry the quotient length metric obtained by collapsing exactly the fat
wedge \(W_{24,T}\) of
\cref{cusp-cone-S24-exact-collapsed-locus}.  Away from that fat wedge the
smash quotient is injective, so no further points are identified.  The
midpoint
\[
 \zeta_{24,T}
 =[(\ell_T/2,\ell_T/2,\ell_T/2,\ell_T/2)]
\]
has distance \(\ell_T/2\) from that fat wedge.  If \(\ell_T>4R\), neither a
nonzero period nor a path through the collapsed fat wedge can shorten a
distance between points in the radius-\(R\) ball of \(C_T^4\).  Conversely,
the collapsed-set distance formula shows that every point of the
radius-\(R\) quotient ball has its unique noncollapsed representative in
that product ball.  Consequently
\begin{equation}
 \rho_{24,T}:
 B_R^{C_T^4}((\ell_T/2,\ell_T/2,\ell_T/2,\ell_T/2))
 \xrightarrow{\ \cong_{\mathrm{isom}}\ }
 B_R^{\mathbb S_{24,T}^4}(\zeta_{24,T})
 \label{cusp-cone-S24-onto-exact-ball}
\end{equation}
is onto as well as isometrically embedded, and its domain is exactly a
Euclidean four-ball because \(\operatorname{sys}(C_T^4)=\ell_T>4R\).  On
the radial source the formula discards only the base angle and then performs
this smash quotient, giving
\cref{cusp-cone-S24-S4-radial-map}.

Both one-copy maps are continuous surjections from compact \(X\) to
Hausdorff spheres, hence quotient maps, and they are based.  Their smash
powers are surjective; compact-to-Hausdorff again makes them quotient maps.
Using the same chosen based diffeomorphism \(X\cong_{\mathrm{diff}}S^6\),
the homeomorphisms
\[
 X^{\wedge2}\cong S^{12},\quad(S^3)^{\wedge2}\cong S^6,
 \qquad
 X^{\wedge4}\cong S^{24},\quad(S^2)^{\wedge4}\cong S^8
\]
prove \cref{cusp-cone-sphere-to-sphere-maps}.  The exact Euclidean-ball
statements prove all three convergences in
\cref{cusp-cone-moving-S4-R4-limits}.
\end{proof}

\begin{corollary}[The two-copy rank-two sphere carrier is null-homotopic]
\label{cusp-cone-S12-sphere-map-null}
After the based identifications in
\cref{cusp-cone-sphere-to-sphere-maps}, the continuous surjective quotient
\[
 \widehat q_3^{\wedge2}:S^{12}\twoheadrightarrow S^6
\]
is based null-homotopic.  In fact, \(g\wedge g:S^{12}\to S^6\) is based
null-homotopic for every based map \(g:S^6\to S^3\).
\end{corollary}

\begin{proof}
Apply the full external-smash calculation in
\cref{cusp-s12-full-carrier-null} to
\(g=\widehat q_3\).  That theorem proves the stronger assertion that
\(x\wedge y:S^{12}\to S^6\) is null-homotopic for every pair of based
maps \(x,y:S^6\to S^3\).
\end{proof}

\begin{corollary}[Every repeated smash of the circle-quotient carrier is null]
\label{cusp-cone-S24-sphere-map-null}
Put \(f=\widehat q_\gamma:S^6\to S^2\) after the chosen based
identification \(X\cong S^6\).  For every \(r\geq2\),
\begin{equation}
 \bigl[f^{\wedge r}\bigr]=0
 \quad\text{in}\quad
 \pi_{6r}(S^{2r}).
 \label{cusp-cone-S24-null-maps}
\end{equation}
In particular,
\begin{equation}
 \bigl[\widehat q_\gamma^{\wedge2}\bigr]
 =0\in\pi_{12}(S^4),
 \qquad
 \bigl[\widehat q_\gamma^{\wedge4}\bigr]
 =0\in\pi_{24}(S^8).
 \label{cusp-cone-S24-null-special-cases}
\end{equation}
Thus the displayed surjective quotient
\(\widehat q_\gamma^{\wedge4}:S^{24}\twoheadrightarrow S^8\) is globally
null-homotopic.  This does not alter either its point-set surjectivity or the
exact intrinsic \(\mathbb R^4\) limit on the moving four-spheres
\(\mathbb S_{24,T}^4\): those are properties of the specified representative
and its radial restrictions, whereas null-homotopy classifies that
representative only up to based homotopy.
\end{corollary}

\begin{proof}
Associativity of the smash product gives the
exact factorization
\begin{equation}
 \begin{tikzcd}[column sep=large]
 S^6\wedge S^6
 \arrow[r,"f\wedge\operatorname{id}_{S^6}"]&
 S^2\wedge S^6
 \arrow[r,"\operatorname{id}_{S^2}\wedge f"]&
 S^2\wedge S^2 .
 \end{tikzcd}
 \label{cusp-cone-S24-smash-square-factorization}
\end{equation}
Under the standard based sphere identifications, the first arrow is the
sixfold suspension of \(f\), up to the harmless twist homeomorphisms dictated
by the left-versus-right suspension convention:
\begin{equation}
 f\wedge\operatorname{id}_{S^6}
 =\Sigma^6f:S^{12}\longrightarrow S^8.
 \label{cusp-cone-S24-sixfold-suspension}
\end{equation}
Freudenthal's suspension theorem is already in its isomorphism range here,
because \(12\leq 2\cdot8-2\), and remains in that range after every further
suspension.  It therefore identifies \(\pi_{12}(S^8)\) with the fourth
stable stem.  The low-stem calculation gives
\begin{equation}
 \pi_{12}(S^8)\cong\pi_4^{\mathrm S}=0
 \quad
 \text{\cite[Corollary~4.24, p.~360]{Hatcher2002};
 \cite[Table~1, p.~113]{IsaksenWangXu2023}}.
 \label{cusp-cone-fourth-stable-stem}
\end{equation}
Hence \(\Sigma^6f\) is based null-homotopic, so the composite in
\cref{cusp-cone-S24-smash-square-factorization}, namely
\(f\wedge f\), is based null-homotopic.  For \(r\geq2\), the canonical
associativity homeomorphisms identify
\begin{equation}
 f^{\wedge r}=(f\wedge f)\wedge f^{\wedge(r-2)},
\end{equation}
with the convention \(f^{\wedge0}=\operatorname{id}_{S^0}\).  Smashing a
based null-homotopy with any based map is a based null-homotopy.  This proves
\cref{cusp-cone-S24-null-maps,cusp-cone-S24-null-special-cases}.
\end{proof}

\begin{corollary}[Literal \(\mathbb R^4\) open cells in the sphere targets]
\label{cusp-cone-literal-R4-open-cells}
Let \(*\) denote the smash basepoint and set
\[
 U_{12,T}:=\mathbb S_{12,T}^4\setminus\{*\},
 \qquad
 U_{24,T}:=\mathbb S_{24,T}^4\setminus\{*\}.
\]
The embedded moving four-spheres restrict to topological embeddings
\[
 U_{12,T}\cong\mathbb R^4\hookrightarrow S^6,
 \qquad
 U_{24,T}\cong\mathbb R^4\hookrightarrow S^8.
\]
More exactly, the quotient-stage maps restrict to homeomorphisms
\begin{align}
 (Q_T\setminus K_T^{(1)})^2
 &\xrightarrow{\ \cong\ }U_{12,T},
 \label{cusp-cone-U12-homeomorphism}\\
 (C_T\setminus\{0\})^4
 &\xrightarrow{\ \cong\ }U_{24,T}.
 \label{cusp-cone-U24-homeomorphism}
\end{align}
These copies of \(\mathbb R^4\) are nonclosed topological open cells: they
are open in their moving four-spheres, not open in the ambient \(S^6\) or
\(S^8\), and their closures in those ambient spheres are
\(\mathbb S_{12,T}^4\) and \(\mathbb S_{24,T}^4\), respectively.  They are
not globally isometric to Euclidean \(\mathbb R^4\) for either the assigned
compact quotient metrics or the restriction of any fixed compatible metric
on the compact sphere targets.  The exact midpoint balls in
\cref{cusp-cone-S12-onto-exact-ball,cusp-cone-S24-onto-exact-ball} give the
precise metric statement: as \(T\to\infty\), every fixed Euclidean
radius-\(R\) ball occurs isometrically and onto the corresponding pointed
ball.
\end{corollary}

\begin{proof}
For a quotient which collapses a closed subset \(K\) to a point, restriction
to the complement gives a homeomorphism
\[
 X\setminus K\xrightarrow{\ \cong\ }(X/K)\setminus\{*\}.
\]
Indeed, the complement of the quotient point is open, its inverse image is
the saturated open set \(X\setminus K\), and restriction of a quotient map
over an open subset of its target is again a quotient map.  The restricted
map is bijective, and a bijective quotient map is a homeomorphism.
Consequently
\[
 Q_T\setminus K_T^{(1)}
 \cong S_T^2\setminus\{*\}.
\]
The marked coordinate one-skeleton \(K^{(1)}\) is the \(1\)-skeleton of the
standard one-vertex CW structure on \(\overline T^2\), with one \(0\)-cell,
two \(1\)-cells, and one \(2\)-cell.  Its complement is therefore the open
\(2\)-cell \((0,1)^2\cong\mathbb R^2\).  Transport by
\(\iota_T\) gives
\[
 Q_T\setminus K_T^{(1)}\cong\mathbb R^2.
\]
Likewise \(C_T\setminus\{0\}\) is an open interval and hence is homeomorphic
to \(\mathbb R\).

The complement of the smash basepoint is the product of the complements of
the factor basepoints.  Therefore
\[
 (S_T^2\wedge S_T^2)\setminus\{*\}
 \cong(S_T^2\setminus\{*\})^2\cong\mathbb R^4,
\]
and
\[
 C_T^{\wedge4}\setminus\{*\}
 \cong(C_T\setminus\{0\})^4\cong\mathbb R^4.
\]
The full moving four-spheres are topologically embedded in their compact
sphere targets, so these restrictions are topological embeddings.  In each
case the omitted smash basepoint lies in the closure of the open cell, which
proves the closure assertions.  Finally, every compatible metric on a
compact target is bounded, as is each compact quotient metric assigned to a
moving four-sphere, whereas Euclidean \(\mathbb R^4\) is unbounded.  Thus no
global isometry is possible; the exact metric content is precisely the
onto-ball calculation already proved.
\end{proof}

\begin{remark}[The moving four-spheres are null-homotopic in the targets]
\label{cusp-cone-moving-S4-null-homotopy}
For a fixed interior reduced-suspension coordinate, the inclusion
\(S_T^2\hookrightarrow S^3\) contracts through based maps by sliding that
coordinate to a suspension endpoint; throughout the homotopy the collapsed
meridian remains at the global basepoint.  Hence
\(S_T^2\wedge S_T^2\hookrightarrow S^6\) is null-homotopic.  The identical
sliding argument contracts \(C_T\hookrightarrow S^2\), and therefore also
contracts \(C_T^{\wedge4}\hookrightarrow S^8\).  Thus the moving metric
four-spheres do not contradict the fixed-target orientation obstruction:
their fundamental classes die in the ambient targets.
\end{remark}

\begin{remark}[Homotopical status of the two sphere-to-sphere maps]
\label{cusp-cone-sphere-map-homotopy-status}
The maps \(S^{12}\to S^6\) and \(S^{24}\to S^8\) induce zero on all
positive-degree integral homology and cohomology groups and have zero
ordinary Hurewicz image.  They have no classical degree because source and
target dimensions differ.  The full four-copy map
\(\widehat q_\gamma^{\wedge4}:S^{24}\to S^8\) is based null-homotopic by
\cref{cusp-cone-S24-sphere-map-null}.  The distinct two-copy map
\(\widehat q_3^{\wedge2}:S^{12}\to S^6\) is also based null-homotopic, by
\cref{cusp-cone-S12-sphere-map-null}.  In both cases the vanishing is a
theorem about the full specified sphere-valued map, not merely its Hurewicz
image or the inclusion of a moving four-sphere. The one-copy class
\([\widehat q_3]\in\pi_6(S^3)\) is now zero by
Theorem~\ref{cusp-null:thm}; the one-copy class
\([\widehat q_\gamma]\in\pi_6(S^2)\) remains open.
\end{remark}

\begin{proposition}[The compact cone-on-\(T^4\) carriers]
\label{cusp-cone-explicit-carriers}
Fix \(T_0\) sufficiently large and put
\(r_0=e^{-2\pi T_0}<\epsilon\).  Define
\[
 \mathcal D_{12}
 =\frac{[0,r_0]\times(\overline T^2)^2}
 {\{0\}\times(\overline T^2)^2\sim d_{12}},
 \qquad
 \mathcal D_{24}
 =\frac{[0,r_0]\times(S^1)^4}
 {\{0\}\times(S^1)^4\sim d_{24}}.
\]
Both spaces are homeomorphic to \(\operatorname{Cone}(T^4)\).  There are
topological embeddings
\begin{align}
 \delta_{12}:\mathcal D_{12}&\hookrightarrow
 \mathcal A\wedge\mathcal A,
 &
 [r,u_1,u_2]&\longmapsto
 [\,[r,1,u_1],[r,1,u_2],],
 \label{cusp-cone-D12-embedding}\\
 \delta_{24}:\mathcal D_{24}&\hookrightarrow
 \mathcal C^{\wedge4},
 &
 [r,\theta_1,\ldots,\theta_4]&\longmapsto
 \bigwedge_{j=1}^4[r,(1,\theta_j)].
 \label{cusp-cone-D24-embedding}
\end{align}
Here \(1\in S^1\) fixes the positive radial base angle.  For every
\(T\geq T_0\), put \(r_T=e^{-2\pi T}\).  The restrictions of the sphere
quotient maps to the synchronized radial fibres are exactly the
cross-sections at height \(r_T\):
\begin{align}
 F_T^2&\xrightarrow{\ \pi_T\times\pi_T\ }
 Q_T^2\cong\{r_T\}\times T^4
 \subset\mathcal D_{12},
 \label{cusp-cone-D12-slice}\\
 F_T^4&\xrightarrow{\ \chi_{\gamma,T}^{\times4}\ }
 C_T^4\cong\{r_T\}\times T^4
 \subset\mathcal D_{24}.
 \label{cusp-cone-D24-slice}
\end{align}
Every one of these embedded four-torus slices is null-homotopic in its
carrier, and hence in the full smash target.

At the cone vertex \(d_\nu\), \(\nu\in\{12,24\}\), the local integral
homology is
\begin{equation}
 H_k(\mathcal D_\nu,
       \mathcal D_\nu\setminus\{d_\nu\};\mathbb Z)
 \cong
 \begin{cases}
  \mathbb Z^4,&k=2,\\
  \mathbb Z^6,&k=3,\\
  \mathbb Z^4,&k=4,\\
  \mathbb Z,&k=5,\\
  0,&\text{otherwise}.
 \end{cases}
 \label{cusp-cone-carrier-local-homology}
\end{equation}
In particular, \(d_\nu\) is not an integral homology-manifold point, with or
without boundary, of any dimension.  Consequently it is not a topological
manifold point, with or without boundary, of any dimension.
\end{proposition}

\begin{proof}
The two displayed domains are cones on
\((\overline T^2)^2\cong T^4\) and \((S^1)^4\cong T^4\).  At \(r=0\),
every factor on the right of
\cref{cusp-cone-D12-embedding,cusp-cone-D24-embedding} is the relevant
smash basepoint, so the formulae descend from the cylinders to the cones.
For \(0<r\leq r_0\), none of the factors is the smash basepoint.  Equality
in the smash quotient is then equality in every factor, and the coordinate
descriptions of the suspension and cusp cone give equality of \(r\) and of
all torus coordinates.  Thus each \(\delta_\nu\) is injective.  Its domain
is compact and its target is Hausdorff, so it is an embedding.

The already proved one-factor restriction formulae, with
\(r_T=e^{-2\pi T}\), are
\[
 q_2(x)=[r_T,1,\iota_T^{-1}(\pi_T(x))]
 \quad(x\in F_T)
\]
and, in the cusp-cone piece of \(\mathcal C\),
\[
 q_1(x)
 =[r_T,(1,\vartheta_T^{-1}(\chi_{\gamma,T}(x)))]
 \quad(x\in F_T).
\]
Applying these identities in two and four factors gives
\cref{cusp-cone-D12-slice,cusp-cone-D24-slice}.  Sliding the cone coordinate
from \(r_T\) to zero contracts each slice to the cone vertex.  Finally,
excision and the long exact sequence of the pair give, for \(L=T^4\),
\[
 H_k(\operatorname{Cone}(L),
     \operatorname{Cone}(L)\setminus\{v\};\mathbb Z)
 \cong\widetilde H_{k-1}(L;\mathbb Z).
\]
The ranks \(4,6,4,1\) of \(H_1,H_2,H_3,H_4\) of \(T^4\) prove
\cref{cusp-cone-carrier-local-homology}.  At an interior point of an
integral homology \(n\)-manifold the local integral homology is \(\mathbb Z\)
in degree \(n\) and zero in every other degree, whereas at a boundary point
of an integral homology manifold with boundary all these local groups
vanish.  The simultaneous nonzero groups in degrees \(2,3,4,5\), including
ranks greater than one, agree with neither pattern for any \(n\).  Every
topological manifold is locally an integral homology manifold, so the final
assertion follows as well.
\end{proof}

\begin{theorem}[Point-set quotient maps and exact homotopy types]
\label{cusp-cone-target-homotopy-types}
The smash maps are continuous surjective quotient maps
\[
 q_2^{\wedge2}:S^{12}\twoheadrightarrow\mathcal A^{\wedge2},
 \qquad
 q_1^{\wedge4}:S^{24}\twoheadrightarrow\mathcal C^{\wedge4}.
\]
Their targets have homotopy types
\begin{align}
 \mathcal A^{\wedge2}
 &\simeq
 \bigvee^9S^4\vee\bigvee^{18}S^5\vee
 \bigvee^{15}S^6\vee\bigvee^6S^7\vee S^8,
 \label{cusp-cone-A-smash-wedge}\\
 \mathcal C^{\wedge4}
 &\simeq
 S^8\vee\bigvee^4S^9\vee\bigvee^6S^{10}
 \vee\bigvee^4S^{11}\vee S^{12}.
 \label{cusp-cone-C-smash-wedge}
\end{align}
Both maps induce zero on every positive-degree integral homology and
cohomology group and have zero ordinary Hurewicz image.  Neither map has a
continuous homotopy section.  These conclusions do not determine whether
the unstable homotopy class of either map is zero.
\end{theorem}

\begin{proof}
Let \(f:(K,k_0)\to(L,\ell_0)\) be a surjective based map between compact
Hausdorff spaces.  Every non-basepoint class of \(L^{\wedge r}\) has a
representative \((y_1,\ldots,y_r)\) with each \(y_i\ne\ell_0\).  Choosing
\(x_i\in f^{-1}(y_i)\) proves surjectivity of \(f^{\wedge r}\); the smash
basepoint is also in its image.  The source smash is compact.  The target
fat wedge is closed, hence the target smash is Hausdorff.  A continuous
surjection from a compact space to a Hausdorff space is a quotient map.
This proves the first assertion for \(q_2\) and \(q_1\).

The suspension splitting
\[
 \Sigma(X_1\times\cdots\times X_n)
 \simeq
 \bigvee_{\varnothing\ne I\subseteq\{1,\ldots,n\}}
 \Sigma\!\left(\bigwedge_{i\in I}X_i\right)
\]
gives
\[
 \mathcal A=\Sigma T^3
 \simeq\bigvee^3S^2\vee\bigvee^3S^3\vee S^4.
\]
Finite smash products distribute over finite wedges.  The ordered degree
sums give multiplicities \(9,18,15,6,1\) in dimensions
\(4,5,6,7,8\), proving \cref{cusp-cone-A-smash-wedge}.

For \(\mathcal C\), write \(x=\rho_1\), \(y=\rho_2\), and
\(h=\eta\).  The pair-of-pants product and the two exact Seifert fillings
give
\[
 \pi_1(\mathcal M)=
 \langle x,y,h\mid[x,h]=[y,h]=1,
 x^3h^{-1}=1,\ y^4h=1\rangle.
\]
The remaining boundary torus is generated by \(h\) and
\(\rho_0=(xy)^{-1}\).  Coning this torus kills their normal closure, and
therefore
\[
 \pi_1(\mathcal C)
 \cong\langle x,y\mid x^3=1,\ y^4=1,\ xy=1\rangle=1,
\]
because \(y=x^{-1}\) and \(\gcd(3,4)=1\).  The previously computed
relative homology gives
\[
 \widetilde H_2(\mathcal C;\mathbb Z)=\mathbb Z,
 \qquad
 \widetilde H_3(\mathcal C;\mathbb Z)=\mathbb Z,
\]
and no other reduced homology.  Since \(\mathcal C\) is a simply connected
finite three-complex, Hurewicz supplies a sphere representing the
\(H_2\)-generator, and the Whitehead exact sequence makes
\(\pi_3(\mathcal C)\to H_3(\mathcal C)\) surjective.  A wedge of
representing maps
\(S^2\vee S^3\to\mathcal C\) is a homology equivalence between simply
connected CW complexes, hence a homotopy equivalence.  Thus
\[
 \mathcal C\simeq S^2\vee S^3.
\]
Taking the fourth smash power yields \cref{cusp-cone-C-smash-wedge} with
the binomial multiplicities \(1,4,6,4,1\).

The first target has no reduced homology above degree eight and the second
has none above degree twelve.  The source spheres have no positive homology
below degrees twelve and twenty-four.  Hence all positive-degree homology
and cohomology homomorphisms induced by the two maps vanish.  Their
Hurewicz images lie in \(H_{12}(\mathcal A^{\wedge2})=0\) and
\(H_{24}(\mathcal C^{\wedge4})=0\), respectively.  A section would split a
nonzero lower homology group of the target into the zero lower homology of
the source sphere, which is impossible.  Higher homotopy groups of the
displayed wedges need not vanish, so none of these arguments decides the
unstable null-homotopy question.
\end{proof}

\begin{theorem}[Exact decompactification spaces with boundary
\(\mathbb R^4\)]
\label{cusp-cone-exact-decompactification-space}
For \(\nu\in\{12,24\}\), set
\[
 A_T^{(12)}=\operatorname{diag}(Y_T,Y_T),
 \qquad
 A_T^{(24)}=\ell_T I_4,
 \qquad s_0=T_0^{-1}.
\]
On \([0,s_0]\times\mathbb R^4\), define an equivalence relation
\(\sim_\nu\) fibrewise over \(s\) by
\begin{align*}
 (0,x)\sim_\nu(0,y)&\iff x=y,\\
 (s,x)\sim_\nu(s,y)&\iff
 x-y\in A_{1/s}^{(\nu)}\mathbb Z^4
 \quad(0<s\leq s_0).
\end{align*}
Let
\[
 \mathfrak E_\nu
 =([0,s_0]\times\mathbb R^4)/\!\sim_\nu,
 \qquad p_\nu:\mathfrak E_\nu\longrightarrow[0,s_0].
\]
Then \(\mathfrak E_\nu\) is a Hausdorff, second-countable topological
five-manifold with boundary, and
\begin{align}
 p_{12}^{-1}(s)&\cong Q_{1/s}^2 &(s>0),&
 &p_{12}^{-1}(0)&\cong\mathbb R^4,
 \label{cusp-cone-E12-fibres}\\
 p_{24}^{-1}(s)&\cong C_{1/s}^4 &(s>0),&
 &p_{24}^{-1}(0)&\cong\mathbb R^4.
 \label{cusp-cone-E24-fibres}
\end{align}
The other boundary component is the compact four-torus over \(s=s_0\).

Give every positive fibre its Euclidean quotient metric and the zero fibre
its Euclidean metric.  For every \(R>0\), there is \(s_R>0\) such that,
for \(0\leq s<s_R\), writing
\(q_{\nu,s}(x)=[s,x]\), the quotient map restricts to an exact isometry
\[
 q_{\nu,s}:B_R^{\mathbb R^4}(0)
 \xrightarrow{\ \cong_{\mathrm{isom}}\ }
 B_R^{p_\nu^{-1}(s)}([s,0])
\]
onto the entire quotient radius-\(R\) ball, not merely onto an embedded
subset of that ball.
Consequently the positive fibres converge to the zero fibre in pointed
local-isometric, pointed smooth, and pointed Gromov--Hausdorff topology.

There are continuous surjective comparison maps which set-theoretically
collapse the zero boundary fibres,
\[
 \begin{tikzcd}[column sep=large]
 \mathbb R^4=p_{12}^{-1}(0)
 \arrow[r,hook]&
 \mathfrak E_{12}
 \arrow[r,"\beta_{12}"]&
 \mathcal D_{12}
 \arrow[r,hook,"\delta_{12}"]&
 \mathcal A^{\wedge2},
 \end{tikzcd}
\]
and
\[
 \begin{tikzcd}[column sep=large]
 \mathbb R^4=p_{24}^{-1}(0)
 \arrow[r,hook]&
 \mathfrak E_{24}
 \arrow[r,"\beta_{24}"]&
 \mathcal D_{24}
 \arrow[r,hook,"\delta_{24}"]&
 \mathcal C^{\wedge4}.
 \end{tikzcd}
\]
For \(s>0\), \(\beta_\nu\) is the specified marked homeomorphism from
\(p_\nu^{-1}(s)\) to the radial torus slice at \(r=e^{-2\pi/s}\).  Its
restriction
\[
 \beta_\nu:\mathfrak E_\nu\setminus p_\nu^{-1}(0)
 \xrightarrow{\ \cong\ }\mathcal D_\nu\setminus\{d_\nu\}
\]
is a homeomorphism, while
\[
 \beta_\nu^{-1}(d_\nu)=p_\nu^{-1}(0)\cong\mathbb R^4.
\]
Neither \(\beta_\nu\) is proper or a quotient map.  It factors through a
continuous bijection
\[
 \overline\beta_\nu:
 \mathfrak E_\nu/p_\nu^{-1}(0)\longrightarrow\mathcal D_\nu,
\]
but this bijection is not a homeomorphism; equivalently, the topology of
\(\mathcal D_\nu\) is not the quotient topology induced by \(\beta_\nu\).
Thus the limiting \(\mathbb R^4\) is an actual boundary fibre of
\(\mathfrak E_\nu\), whereas the continuous comparison map sends that entire
fibre to the cone vertex.
\end{theorem}

\begin{proof}
We give a common proof.  Write
\[
 L_s=A_{1/s}^{(\nu)}\mathbb Z^4\quad(s>0),
 \qquad L_0=\{0\}.
\]
The matrices vary continuously for \(s>0\), and
\[
 \operatorname{sys}(L_s)\longrightarrow\infty
 \qquad(s\downarrow0)
\]
by \(\operatorname{sys}(Y_T\mathbb Z^2)\geq T-C\) in the
\(S^{12}\) case and \(\operatorname{sys}(\ell_T\mathbb Z^4)=\ell_T\)
in the \(S^{24}\) case.

The equivalence relation is closed.  Indeed, suppose
\((s_j,x_j)\sim_\nu(s_j,y_j)\) and
\((s_j,x_j,y_j)\to(s,x,y)\).  If \(s>0\), write
\(y_j-x_j=A_{1/s_j}^{(\nu)}n_j\).  On a compact neighborhood of \(s\),
the inverse matrices are uniformly bounded, so boundedness of
\(y_j-x_j\) forces \(n_j\) to lie in a finite set.  Passing to a
subsequence makes \(n_j\) constant and proves \(x\sim_\nu y\).  If
\(s=0\), any nonzero \(n_j\) would give
\[
 \|y_j-x_j\|\geq\operatorname{sys}(L_{s_j})\longrightarrow\infty,
\]
contrary to convergence.  Hence \(x=y\).

The quotient projection \(q_\nu:[0,s_0]\times\mathbb R^4\to
\mathfrak E_\nu\) is open.  If \(U\) is open, its saturation is the union
of \(U\) with the images, for \(n\in\mathbb Z^4\), of
\[
 U\cap((0,s_0]\times\mathbb R^4)
 \quad\text{under}\quad
 (s,x)\longmapsto(s,x+A_{1/s}^{(\nu)}n),
\]
and every one of these images is open.  A closed equivalence relation with
an open quotient map has Hausdorff quotient: two inequivalent points admit
a product neighborhood disjoint from the relation, and the saturations of
the two factor neighborhoods are disjoint open sets.  Since the source is
second countable and the quotient map is open, images of a countable base
form a countable base of the quotient.

Fix \(x_0\in\mathbb R^4\).  Choose \(a>0\) and \(s_1>0\) such that
\(\operatorname{sys}(L_s)>2a\) for \(0<s<s_1\).  The quotient projection
is one-to-one on
\([0,s_1)\times B_a^\circ(x_0)\), where
\(B_a^\circ(x_0)=\{x:\|x-x_0\|<a\}\), and its restriction is an open continuous
bijection onto an open subset of \(\mathfrak E_\nu\); it is therefore a
homeomorphism.  These are half-space charts at \(s=0\).  The usual evenly
covered small balls give product charts for \(s>0\), and the same argument
at \(s=s_0\) gives the second boundary component.  This proves the
manifold assertion and the fibre identifications.

For the metric assertion, choose \(s_R\) so that
\(\operatorname{sys}(L_s)>4R\) for \(s<s_R\).  If
\(x,y\in B_R(0)\) and \(0\ne\lambda\in L_s\), then
\[
 \|x-y+\lambda\|
 \geq\|\lambda\|-\|x-y\|
 >4R-2R\geq\|x-y\|.
\]
Thus no translated representative shortens the Euclidean distance on that
ball, so \(q_{\nu,s}\) is an exact isometric embedding there.  It is also
onto the full quotient ball.  Indeed, if \(\xi=[s,x]\) satisfies
\(d([s,0],\xi)\leq R\), discreteness of \(L_s\) makes the infimum in
\[
 d([s,0],[s,x])=\min_{\lambda\in L_s}\|x+\lambda\|
\]
attained.  For a minimizing \(\lambda_0\), the representative
\(x+\lambda_0\) lies in \(B_R^{\mathbb R^4}(0)\) and maps to \(\xi\).  At
\(s=0\), \(q_{\nu,0}\) is the identity.  This proves the asserted onto
isometry, and since it holds for every fixed \(R\), it proves all three
pointed convergences.

It remains to define these comparison maps, which set-theoretically collapse
the zero fibres.  For \(s>0\), put
\(T=1/s\) and \(r=e^{-2\pi T}\).  In the \(S^{12}\) case, write
\(x=(x_1,x_2)\in\mathbb R^2\oplus\mathbb R^2\) and set
\[
 \beta_{12}([s,x])
 =\bigl[r,
 \iota_T^{-1}([x_1]_{Y_T\mathbb Z^2}),
 \iota_T^{-1}([x_2]_{Y_T\mathbb Z^2})\bigr].
\]
In the \(S^{24}\) case, write
\(x=(x_1,\ldots,x_4)\) and set
\[
 \beta_{24}([s,x])
 =\bigl[r,
 [x_1/\ell_T],\ldots,[x_4/\ell_T]\bigr].
\]
These formulae are well defined on the positive fibres and are continuous
there.  Define \(\beta_\nu([0,x])=d_\nu\) for every \(x\in\mathbb R^4\).
Every neighborhood of \(d_\nu\) in the cone contains all cross-sections
with sufficiently small \(r\).  Since \(e^{-2\pi/s}\to0\) as
\(s\downarrow0\), continuity at the zero fibre is uniform in \(x\).
Surjectivity follows because each positive fibre maps onto its complete
cross-section and the zero fibre maps onto the vertex.  The restriction away
from \(p_\nu^{-1}(0)\) is a homeomorphism because
\(s\mapsto e^{-2\pi/s}\) is a homeomorphism from \((0,s_0]\) to
\((0,r_0]\), and the displayed marked-torus identifications are fibrewise
homeomorphisms varying continuously with \(s\).  The vertex fibre is exactly
\(p_\nu^{-1}(0)\), so \(\beta_\nu\) is not proper.

It is also not a quotient map.  Choose
\(a\in\mathbb R^4\setminus\mathbb Z^4\), for example
\(a=(1/2,1/2,1/2,1/2)\), choose \(s_j\downarrow0\), and put
\[
 z_j=[s_j,A_{1/s_j}^{(\nu)}a]\in\mathfrak E_\nu.
\]
For \(\nu=12\),
\[
 \inf_{n\in\mathbb Z^4}
 \|A_{1/s_j}^{(12)}(a+n)\|
 \geq(1/s_j-C)\operatorname{dist}(a,\mathbb Z^4)
 \longrightarrow\infty,
\]
and for \(\nu=24\) the same infimum equals
\[
 \ell_{1/s_j}\operatorname{dist}(a,\mathbb Z^4)
 \longrightarrow\infty.
\]
More explicitly, for every fixed \(x_0\in\mathbb R^4\),
\[
 \inf_{n\in\mathbb Z^4}
 \|A_{1/s_j}^{(\nu)}(a+n)-x_0\|
 \geq
 \inf_{n\in\mathbb Z^4}\|A_{1/s_j}^{(\nu)}(a+n)\|-\|x_0\|
 \longrightarrow\infty.
\]
Hence \(\{z_j:j\geq1\}\) is closed in \(\mathfrak E_\nu\): it has no
accumulation in a positive fibre and the last estimate shows that it
eventually misses every half-space chart about every point of the zero
fibre.  But
\[
 \beta_\nu(z_j)=[e^{-2\pi/s_j},a]\longrightarrow d_\nu,
\]
so \(\{\beta_\nu(z_j):j\geq1\}\) is not closed in \(\mathcal D_\nu\).
Since \(\beta_\nu\) is injective off the zero fibre and the latter maps to
\(d_\nu\), the inverse image of this nonclosed set is the closed set
\(\{z_j:j\geq1\}\).  Thus \(\beta_\nu\) is not a quotient map.  This proves
the two diagrams and their exact topological scope.  Since \(\beta_\nu\) is
constant exactly on \(p_\nu^{-1}(0)\) and otherwise injective, it induces the
continuous bijection \(\overline\beta_\nu\) displayed in the theorem.  Were
that bijection a homeomorphism, \(\beta_\nu\) would be the composite of the
quotient projection with a homeomorphism and hence a quotient map, contrary
to the counterexample.  Therefore \(\overline\beta_\nu\) is not a
homeomorphism.
\end{proof}

\begin{corollary}[Sharp global meaning of the Euclidean limit]
\label{cusp-cone-sharp-global-R4-scope}
The native routes constructed above exhibit the underlying space
\(\mathbb R^4\) in the following three precise senses:
\begin{enumerate}
\item as each of the nonclosed topologically embedded open cells
\(U_{12,T}\subset S^6\) and \(U_{24,T}\subset S^8\);
\item as the pointed intrinsic Euclidean limit of the four-torus slices in the
global quotient carriers of \(S^{12}\) and \(S^{24}\); and
\item as the Euclidean decompactification boundary fibre of
\(\mathfrak E_{12}\) and \(\mathfrak E_{24}\).
\end{enumerate}
It does not produce \(\mathbb R^4\) as the whole compact sphere, as the
whole compact conical target, or as a continuous quotient target of either
compact space.  Nor does the unrescaled restriction of any fixed compatible
ambient target metric have \(\mathbb R^4\) as its slice limit, and the open
cells in item \(1\) are not globally Euclidean-isometric in such a metric.
\end{corollary}

\begin{proof}
The three positive assertions are
\cref{cusp-cone-literal-R4-open-cells,cusp-cone-explicit-carriers,%
cusp-cone-exact-decompactification-space}.
For the negative assertions, a continuous image of a compact sphere or
compact conical target is compact, so it cannot be all of \(\mathbb R^4\).
The local homology calculation
\cref{cusp-cone-carrier-local-homology} shows that the carrier vertex is not
an integral homology-manifold point, and hence not a manifold point, of any
dimension, with or without boundary.  Finally, in every fixed compatible
metric on a compact cone target, all sufficiently deep radial slices lie
uniformly in an arbitrarily small neighborhood of the vertex, so their ambient
diameters tend to zero.  Their intrinsic diameters tend to infinity because
their intrinsic systoles do.  Hence without a \(T\)-dependent rescaling no
fixed ambient metric realizes the intrinsic Euclidean limit on those slices.
The fixed cusp-adapted blow-ups are treated separately in
\cref{sphere-metric-fixed-realization}.  The comparison diagrams record the
unrescaled distinction exactly: they collapse the new \(\mathbb R^4\)
boundary rather than identify it with the compact target.
\end{proof}

\paragraph{Scope.}
The constructions above use only the original, unscaled fibre metrics and
the quotient maps derived from the native \(S^6\) cusp.  They do not use the
auxiliary finite-cover unspooling.  The conical carriers depend on the
chosen cusp coordinate, positive radial ray, and based identification
\(X\cong S^6\).  The spaces \(\mathfrak E_\nu\) additionally remember the
marked Euclidean quotient metrics.  None of these data defines a compact
complex structure on \(S^{12}\) or \(S^{24}\), a four-dimensional physical
spacetime, or an interacting Yang--Mills Hamiltonian.
